% GENERATED FILE. Edit canonical lessons, then run npm run generate:books.
% canonical-source-hash: fnv1a64:a23c5556e9bf7004
% canonical-lessons: TA-C550-astivaram, TA-C550-cenkal, TA-C550-ciment, TA-C550-caram, TA-C550-mestiri

\chapter{Building a House}
\label{ch:ta-building-a-house}
\hlchaptermodality{550}

\begin{chapteropening}
\textbf{By the end of this chapter:} \emph{I can talk about bricks, cement and building work.}

Complete the last lesson of chapter 550: I can talk about bricks, cement and building work.
\end{chapteropening}

\section[\texorpdfstring{\ta{அஸ்திவாரம்} (astivāram)}{அஸ்திவாரம் (astivāram)}]{\ta{அஸ்திவாரம்} (astivāram) — a foundation}

\label{lesson:TA-C550-astivaram}



\begin{quote}
Before the new one: say the Tamil for readiness, then the Tamil for to slip away.
\end{quote}

\subsection*{You'll want to know: \ta{அஸ்திவாரம்}}
\textbf{\ta{அஸ்திவாரம்}} — \emph{astivāram} — \textquotedblleft{}a foundation\textquotedblright{}.

\begin{quote}
\ta{அஸ்திவாரம்} \ta{பலமாக} \ta{இல்லாவிட்டால்} \ta{வீடு} \ta{நிற்காது}.
\end{quote}

\begin{quote}
\emph{astivāram palamāka illāviṭṭāl vīṭu niṟkātu.}
\end{quote}

\begin{quote}
If the foundation is not strong, the house will not stand.
\end{quote}

Also of ideas: \ta{நம்பிக்கைதான்} \ta{உறவின்} \ta{அஸ்திவாரம்}, trust is the foundation of a relationship.

\begin{cousinweb}
The origin of \textbf{\ta{அஸ்திவாரம்}} (\emph{astivāram}) is not settled: no dictionary we can check gives its source.
\end{cousinweb}

\begin{grammarlens}[title={What you've built}]
One more word for saying what happens if.
\end{grammarlens}

\subsection*{Your turn}
\begin{itemize}
\raggedright
  \item \emph{Say it:} \emph{astivāram}
  \item \emph{Say it:} \emph{astivāram}, once more
  \item \emph{Say it:} \emph{astivāram palamāka illāviṭṭāl vīṭu niṟkātu.} — If the foundation is not strong, the house will not stand.
  \item \emph{Recall:} say the Tamil for readiness, then the Tamil for to slip away, then say \emph{astivāram} again
\end{itemize}

\begin{tcolorbox}[breakable,colback=teal!4,colframe=teal!35!black,title={Before you move on}]
What does \ta{அஸ்திவாரம்} mean? (\textquotedblleft{}A foundation\textquotedblright{}.) Say it once more.
\end{tcolorbox}
\section[\texorpdfstring{\ta{செங்கல்} (ceṅkal)}{செங்கல் (ceṅkal)}]{\ta{செங்கல்} (ceṅkal) — a brick}

\label{lesson:TA-C550-cenkal}



\begin{quote}
Before the new one: say the Tamil for to slip away, then the Tamil for a foundation.
\end{quote}

\subsection*{You'll want to know: \ta{செங்கல்}}
\textbf{\ta{செங்கல்}} — \emph{ceṅkal} — \textquotedblleft{}a brick\textquotedblright{}.

\begin{quote}
\ta{லாரியில்} \ta{ஆயிரம்} \ta{செங்கற்கள்} \ta{வந்தன}.
\end{quote}

\begin{quote}
\emph{lāriyil āyiram ceṅkaṟkaḷ vantaṉa.}
\end{quote}

\begin{quote}
A thousand bricks came on the lorry.
\end{quote}

Plural: \ta{செங்கற்கள்}, \emph{ceṅkaṟkaḷ}.

Before -\ta{கள்} the final -\ta{ல்} becomes -\ta{ற்}: \ta{செங்கற்கள்}, bricks.

\begin{cousinweb}
\textbf{\ta{செங்கல்}} (\emph{ceṅkal}) is a native compound: \ta{செம்} (\emph{cem}), \textquotedblleft{}red\textquotedblright{} (DEDR 1931), and \ta{கல்} (\emph{kal}), \textquotedblleft{}a stone\textquotedblright{} (DEDR 1298): a red stone.
\end{cousinweb}

\begin{grammarlens}[title={What you've built}]
One more word for saying what happens if.
\end{grammarlens}

\subsection*{Your turn}
\begin{itemize}
\raggedright
  \item \emph{Say it:} \emph{ceṅkal}
  \item \emph{Say it:} \emph{ceṅkal}, once more
  \item \emph{Say it:} \emph{lāriyil āyiram ceṅkaṟkaḷ vantaṉa.} — A thousand bricks came on the lorry.
  \item \emph{Recall:} say the Tamil for to slip away, then the Tamil for a foundation, then say \emph{ceṅkal} again
\end{itemize}

\begin{tcolorbox}[breakable,colback=teal!4,colframe=teal!35!black,title={Before you move on}]
What does \ta{செங்கல்} mean? (\textquotedblleft{}A brick\textquotedblright{}.) Say it once more.
\end{tcolorbox}
\section[\texorpdfstring{\ta{சிமெண்ட்} (cimeṇṭ)}{சிமெண்ட் (cimeṇṭ)}]{\ta{சிமெண்ட்} (cimeṇṭ) — cement}

\label{lesson:TA-C550-ciment}



\begin{quote}
Before the new one: say the Tamil for a foundation, then the Tamil for a brick.
\end{quote}

\subsection*{You'll want to know: \ta{சிமெண்ட்}}
\textbf{\ta{சிமெண்ட்}} — \emph{cimeṇṭ} — \textquotedblleft{}cement\textquotedblright{}.

\begin{quote}
\ta{சுவர்} \ta{கட்ட} \ta{இன்னும்} \ta{பத்து} \ta{மூட்டை} \ta{சிமெண்ட்} \ta{வேண்டும்}.
\end{quote}

\begin{quote}
\emph{cuvar kaṭṭa iṉṉum pattu mūṭṭai cimeṇṭ vēṇṭum.}
\end{quote}

\begin{quote}
We need ten more sacks of cement to build the wall.
\end{quote}

Written with the final \ta{ட்} of the English word.

\begin{cousinweb}
\textbf{\ta{சிமெண்ட்}} (\emph{cimeṇṭ}) is a loan from English \emph{cement}.
\end{cousinweb}

\begin{grammarlens}[title={What you've built}]
One more word for saying what happens if.
\end{grammarlens}

\subsection*{Your turn}
\begin{itemize}
\raggedright
  \item \emph{Say it:} \emph{cimeṇṭ}
  \item \emph{Say it:} \emph{cimeṇṭ}, once more
  \item \emph{Say it:} \emph{cuvar kaṭṭa iṉṉum pattu mūṭṭai cimeṇṭ vēṇṭum.} — We need ten more sacks of cement to build the wall.
  \item \emph{Recall:} say the Tamil for a foundation, then the Tamil for a brick, then say \emph{cimeṇṭ} again
\end{itemize}

\begin{tcolorbox}[breakable,colback=teal!4,colframe=teal!35!black,title={Before you move on}]
What does \ta{சிமெண்ட்} mean? (\textquotedblleft{}Cement\textquotedblright{}.) Say it once more.
\end{tcolorbox}
\section[\texorpdfstring{\ta{சாரம்} (cāram)}{சாரம் (cāram)}]{\ta{சாரம்} (cāram) — scaffolding}

\label{lesson:TA-C550-caram}



\begin{quote}
Before the new one: say the Tamil for a brick, then the Tamil for cement.
\end{quote}

\subsection*{You'll want to know: \ta{சாரம்}}
\textbf{\ta{சாரம்}} — \emph{cāram} — \textquotedblleft{}scaffolding\textquotedblright{}.

\begin{quote}
\ta{தொழிலாளர்கள்} \ta{சாரத்தின்} \ta{மேல்} \ta{நின்று} \ta{வேலை} \ta{செய்தார்கள்}.
\end{quote}

\begin{quote}
\emph{toḻilāḷarkaḷ cārattiṉ mēl niṉṟu vēlai ceytārkaḷ.}
\end{quote}

\begin{quote}
The workers stood on the scaffolding and worked.
\end{quote}

A -\ta{அம்} noun: \ta{சாரத்தில்}, \ta{சாரத்தின்} \ta{மேல்}.

\begin{cousinweb}
The origin of \textbf{\ta{சாரம்}} (\emph{cāram}) in this sense is not settled; it sounds the same as \ta{சாரம்}, \textquotedblleft{}essence\textquotedblright{}, a different word.
\end{cousinweb}

\begin{grammarlens}[title={What you've built}]
One more word for saying what happens if.
\end{grammarlens}

\subsection*{Your turn}
\begin{itemize}
\raggedright
  \item \emph{Say it:} \emph{cāram}
  \item \emph{Say it:} \emph{cāram}, once more
  \item \emph{Say it:} \emph{toḻilāḷarkaḷ cārattiṉ mēl niṉṟu vēlai ceytārkaḷ.} — The workers stood on the scaffolding and worked.
  \item \emph{Recall:} say the Tamil for a brick, then the Tamil for cement, then say \emph{cāram} again
\end{itemize}

\begin{tcolorbox}[breakable,colback=teal!4,colframe=teal!35!black,title={Before you move on}]
What does \ta{சாரம்} mean? (\textquotedblleft{}Scaffolding\textquotedblright{}.) Say it once more.
\end{tcolorbox}
\section[\texorpdfstring{\ta{மேஸ்திரி} (mēstiri)}{மேஸ்திரி (mēstiri)}]{\ta{மேஸ்திரி} (mēstiri) — a foreman; a head mason}

\label{lesson:TA-C550-mestiri}



\begin{quote}
Before the new one: say the Tamil for cement, then the Tamil for scaffolding.
\end{quote}

\subsection*{You'll want to know: \ta{மேஸ்திரி}}
\textbf{\ta{மேஸ்திரி}} — \emph{mēstiri} — \textquotedblleft{}a foreman; a head mason\textquotedblright{}.

\begin{quote}
\ta{மேஸ்திரி} \ta{தொழிலாளர்களுக்கு} \ta{வேலையைப்} \ta{பிரித்துக்} \ta{கொடுத்தார்}.
\end{quote}

\begin{quote}
\emph{mēstiri toḻilāḷarkaḷukku vēlaiyaip pirittuk koṭuttār.}
\end{quote}

\begin{quote}
The foreman shared out the work among the workers.
\end{quote}

A head of a building gang; compare \ta{கொத்தனார்} (\emph{kottaṉār}), a mason.

\begin{cousinweb}
\textbf{\ta{மேஸ்திரி}} (\emph{mēstiri}) is a loan from Portuguese \emph{mestre}, \textquotedblleft{}a master\textquotedblright{}, which also reached other Indian languages.
\end{cousinweb}

\begin{grammarlens}[title={What you've built}]
One more word for saying what happens if.
\end{grammarlens}

\subsection*{Your turn}
\begin{itemize}
\raggedright
  \item \emph{Say it:} \emph{mēstiri}
  \item \emph{Say it:} \emph{mēstiri}, once more
  \item \emph{Say it:} \emph{mēstiri toḻilāḷarkaḷukku vēlaiyaip pirittuk koṭuttār.} — The foreman shared out the work among the workers.
  \item \emph{Recall:} say the Tamil for cement, then the Tamil for scaffolding, then say \emph{mēstiri} again
\end{itemize}

\begin{tcolorbox}[breakable,colback=teal!4,colframe=teal!35!black,title={Before you move on}]
What does \ta{மேஸ்திரி} mean? (\textquotedblleft{}A foreman; a head mason\textquotedblright{}.) Say it once more.
\end{tcolorbox}
